% Shared bilingual method view: tokens -> backbone -> vectors -> head -> scores.
% Colored blocks denote object types, not measured activations or probabilities.
% Numeric scores illustrate the decision rules; they are not model predictions.
\documentclass[tikz,border=2mm]{standalone}
\usepackage[UTF8,fontset=none]{ctex}
\setmainfont{DejaVu Sans}
\setsansfont{DejaVu Sans}
\setCJKmainfont{Noto Sans CJK SC}
\setCJKsansfont{Noto Sans CJK SC}
\usepackage{amsmath,amssymb}
\usepackage{tikz}
\usetikzlibrary{arrows.meta,calc,backgrounds}
\usepackage{hyperref}
\hypersetup{pdfauthor={shichao},pdftitle={Metis candidate scoring method},
  pdfsubject={Illustrated Base plus ScoreHead forward and LoRA training},
  pdfkeywords={Metis, candidate scoring, ScoreHead},pdfcreator={TeX}}
\definecolor{Ink}{HTML}{273648}
\definecolor{Muted}{HTML}{607086}
\definecolor{Data}{HTML}{408AA4}
\definecolor{Hidden}{HTML}{8C73B5}
\definecolor{Weight}{HTML}{D89548}
\definecolor{Score}{HTML}{C96E78}
\definecolor{Target}{HTML}{7E8B9B}
\definecolor{Chosen}{HTML}{388878}
\newcommand{\lang}[2]{\ifMetisChinese#2\else#1\fi}
\newcommand{\Small}{\fontsize{8.5}{10.5}\selectfont}
\newcommand{\Body}{\fontsize{10}{12}\selectfont}
\newcommand{\Heading}{\fontsize{11}{13}\selectfont\bfseries}
\newcommand{\Cell}[5]{%
  \path[draw=#4!45,fill=#4!#5,rounded corners=.5pt,line width=.3pt]
    (#1,#2) rectangle ++(#3,#3);}
\newcommand{\Swatch}[3]{%
  \path[fill=#1!65,draw=#1!65,rounded corners=1pt] (#2,.57) rectangle ++(.19,.19);
  \node[anchor=west,font=\Small] at (#2+.29,.665) {#3};}
\begin{document}
\sffamily
\begin{tikzpicture}[x=1cm,y=1cm,text=Ink,font=\Body,
  flow/.style={-{Latex[length=2mm,width=1.35mm]},draw=Ink,line width=.8pt},
  gradient/.style={-{Latex[length=2mm,width=1.35mm]},draw=Score,line width=.85pt,dashed},
  block/.style={draw=Ink!25,fill=white,rounded corners=2pt,line width=.55pt,align=center},
  note/.style={font=\Small,text=Muted,align=center}]
\path[use as bounding box] (0,0) rectangle (23,10.6);
\begin{scope}[on background layer]
  \fill[white] (0,0) rectangle (23,10.6);
  \path[fill=Data!6,rounded corners=5pt] (.1,5.75) rectangle (5.3,9.15);
  \path[fill=Hidden!6,rounded corners=5pt] (5.65,5.75) rectangle (16.15,9.15);
  \path[fill=Score!6,rounded corners=5pt] (16.5,5.75) rectangle (22.9,9.15);
  \path[fill=Score!4,rounded corners=5pt] (.1,1.3) rectangle (22.9,4.55);
\end{scope}

\node[anchor=west,font=\fontsize{15}{18}\selectfont\bfseries] at (.1,10.25)
  {\lang{Metis: one score for each candidate}{Metis：把每个候选变成一个分数}};
\node[anchor=west,note] at (.1,9.65)
  {\lang{Shared backbone and head; candidates may be processed in batches.}{所有候选共用基座和评分头，可按候选分批计算。}};
\node[anchor=east,note] at (22.9,9.65)
  {Qwen3 Base + ScoreHead};

% Main path: colored areas show the three stages; rows preserve candidate order.
\node[anchor=west,font=\Heading,text=Data!85!Ink] at (.4,8.81)
  {\lang{1. Candidate inputs}{1. 逐候选组织输入}};
\node[note] at (1.75,8.33) {\lang{Shared}{公共输入}};
\node[note] at (3.34,8.33) {\lang{Text}{候选正文}};
\node[note] at (4.58,8.33) {\lang{Suffix}{末尾}};
\foreach \y/\i in {7.85/1,7.10/2,6.35/3}{
  \node[font=\Small,text=Muted] at (.43,\y) {$c_\i$};
  \node[draw=Data!45,fill=Data!15,rounded corners=1pt,
        minimum width=1.28cm,minimum height=.46cm,font=\Small] at (1.7,\y) {$q,u$};
  \node[draw=Data!65,fill=Data!40,rounded corners=1pt,
        minimum width=1.63cm,minimum height=.46cm,font=\Small] at (3.3,\y)
    {\ifnum\i=1\lang{Rules}{退票规则}\fi
     \ifnum\i=2\lang{Weather}{天气预报}\fi
     \ifnum\i=3\lang{Fees}{退票手续费}\fi};
  \node[draw=Data!80,fill=Data!65,text=white,rounded corners=1pt,
        minimum width=.38cm,minimum height=.46cm,font=\Small] at (4.55,\y) {R};
}
\node[note] at (2.72,5.24)
  {\lang{$q$: refund information\\$u$: context; R: last suffix token}{$q$：退票相关（示例问题）\\$u$：可选上下文；R：末尾读出}};

\node[anchor=west,font=\Heading,text=Hidden!85!Ink] at (5.95,8.81)
  {\lang{2. Encode and score}{2. 取出表示，再映射为分数}};
\foreach \off in {.16,.08,0}{
  \path[draw=Hidden!45,fill=white,rounded corners=2pt]
    (6.05+\off,6.22+\off) rectangle (8.78+\off,8.24+\off);
}
\node[font=\Heading] at (7.42,7.68) {Qwen3 Base};
\node[note] at (7.42,7.15) {\lang{Shared weights}{每行使用同一个模型}};
\node[draw=Weight!55,fill=Weight!17,rounded corners=2pt,
      minimum width=1.92cm,minimum height=.39cm,font=\Small] at (7.42,6.63)
  {LoRA $\Delta\theta$};
\draw[flow] (4.94,7.13) -- (5.91,7.13);

% One D-dimensional vector from each candidate's last valid suffix position.
\node[note] at (10.28,8.34) {\lang{Readout vectors}{末尾位置向量}};
\foreach \y/\i in {7.85/1,7.10/2,6.35/3}{
  \node[font=\Small] at (9.44,\y) {$h_\i$};
  \foreach \j in {0,...,3}{
    \pgfmathtruncatemacro{\tone}{30+mod(\i+\j,3)*22}
    \Cell{9.79+\j*.31}{\y-.13}{.26}{Hidden}{\tone}
  }
}
\draw[flow] (8.99,7.13) -- (9.25,7.13);
\node[note] at (10.32,5.99) {$h_i\in\mathbb{R}^{D}$};

% Small weight matrices make the head visible without expanding every operator.
\node[note] at (13.64,8.34) {\lang{Shared MLP head}{共用的 MLP 评分头}};
\foreach \r in {0,...,3}{\foreach \c in {0,...,2}{
  \pgfmathtruncatemacro{\tone}{25+mod(\r+\c,3)*22}
  \Cell{12.03+\c*.22}{6.76+\r*.22}{.18}{Weight}{\tone}
}}
\node[block,font=\Small,minimum height=.49cm,minimum width=1.05cm]
  at (13.62,7.13) {GELU};
\foreach \r in {0,...,2}{\Cell{14.63}{6.85+\r*.22}{.18}{Weight}{65}}
\node[note] at (12.35,6.35) {$W_1,b_1$};
\node[note] at (14.73,6.35) {$W_2,b_2$};
\draw[flow] (11.13,7.13) -- (11.84,7.13);
\draw[flow] (12.76,7.13) -- (13.02,7.13);
\draw[flow] (14.22,7.13) -- (14.44,7.13);
\node[note] at (13.58,5.99) {$s_i=\operatorname{MLP}_{\phi}(h_i)$};

\node[anchor=west,font=\Heading,text=Score!85!Ink] at (16.8,8.81)
  {\lang{3. Task decisions}{3. 根据分数作决策}};
\node[note] at (17.99,8.33) {\lang{Raw scores}{原始分数}};
\draw[Ink!20,line width=.45pt] (17.8,6.04) -- (17.8,8.14);
\foreach \y/\i/\value/\width in {7.85/1/2.1/.84,7.10/2/-0.8/-.32,6.35/3/1.4/.56}{
  \node[font=\Small,text=Muted] at (16.99,\y) {$c_\i$};
  \path[draw=Score!70,fill=Score!65,rounded corners=.6pt]
    (17.8,\y-.13) rectangle (17.8+\width,\y+.13);
  \node[font=\Small,anchor=east] at (19.25,\y) {$\value$};
}
\draw[flow] (14.99,7.13) -- (16.30,7.13);
\draw[flow] (19.43,7.13) -- (19.89,7.13);
\node[note] at (21.32,8.33) {\lang{Task rule}{按任务选择规则}};
\node[block,draw=Chosen!40,fill=Chosen!8,minimum width=2.45cm,
      minimum height=.49cm,font=\Small] at (21.33,7.85)
  {\lang{Rank}{排序} $c_1,c_3,c_2$};
\node[block,draw=Chosen!40,fill=Chosen!8,minimum width=2.45cm,
      minimum height=.49cm,font=\Small] at (21.33,7.10)
  {\lang{Choose}{单选} $c_1$};
\node[block,draw=Chosen!40,fill=Chosen!8,minimum width=2.45cm,
      minimum height=.49cm,font=\Small] at (21.33,6.35)
  {\lang{Multi}{多选} $\{c_1,c_3\}$};
\node[note] at (19.8,5.24)
  {\lang{Illustrative scores\\Multi-label example: $s_i\geq0$}{分数为示意值\\多选示例的 logit 阈值为 $0$}};

% Training is a separate, short addition to the same forward path.
\node[anchor=west,font=\Heading,text=Score!85!Ink] at (.4,4.18)
  {\lang{Training: loss and parameter update}{训练只多两步：计算损失，更新参数}};
\node[note] at (2.37,3.5) {\lang{Scores}{上方得到的分数}};
\node[note] at (5.65,3.5) {\lang{Labels (ranking)}{监督标签（排序例）}};
\foreach \y/\score/\label in {3.0/2.1/1,2.58/-0.8/0,2.16/1.4/1}{
  \node[draw=Score!40,fill=Score!22,minimum width=1.18cm,
    minimum height=.32cm,font=\Small] at (2.37,\y) {$\score$};
  \node[draw=Target!40,fill=Target!20,minimum width=1.18cm,
    minimum height=.32cm,font=\Small] at (5.65,\y) {$\label$};
}
\node[block,draw=Score!40,fill=white,minimum width=3.0cm,
  minimum height=1.35cm] (loss) at (10.8,2.71)
  {\lang{Task loss}{任务损失}\\[5pt]$\mathcal{L}(\boldsymbol{s},\boldsymbol{y})$};
\draw[flow] (3.05,2.85) -- (3.65,2.85) -- (3.65,3.8) -- (10.8,3.8) -- (loss.north);
\draw[flow] (6.35,2.43) -- (loss.west |- 0,2.43);
\node[block,draw=Weight!50,fill=Weight!8,minimum width=5.22cm,
  minimum height=1.35cm] (update) at (17.95,2.71)
  {\lang{Update LoRA + head}{更新 LoRA 与评分头}\\[5pt]
   $\Delta\theta\,,\quad \phi=(W_1,b_1,W_2,b_2)$};
\draw[gradient] (loss.east) -- node[note,above,text=Score]
  {\lang{Backpropagation}{反向传播}} (update.west);
\node[note,anchor=west] at (.4,1.63)
  {\lang{LoRA example: base weights $\theta_0$ stay fixed; labels enter only the loss. The toolkit also supports full fine-tuning.}{LoRA 示例：基座原权重 $\theta_0$ 冻结；标签只进入损失。工具包另支持全量微调。}};

\Swatch{Data}{.4}{\lang{Input tokens}{输入 token}}
\Swatch{Hidden}{4.3}{\lang{Hidden vectors}{隐藏向量}}
\Swatch{Weight}{8.3}{\lang{Head weights}{评分头参数}}
\Swatch{Score}{12.5}{\lang{Scores / loss}{分数 / 损失}}
\Swatch{Target}{17.2}{\lang{Labels}{监督标签}}
\node[anchor=west,font=\fontsize{7.5}{9}\selectfont,text=Muted] at (.4,.13)
  {\lang{IDs follow the scores; no text generation. Colored cells illustrate object types, not measured values.}{候选 ID 随分数保留，不生成文本。色块表示对象类型，不表示实际激活值。}};
\end{tikzpicture}
\end{document}
